\section{Varying positive backgrounds}\label{sec:geometric-normalization-stability}

A two-sided inequality between two Hermitian forms compares their normalization scalars uniformly in the logarithmic density. This turns geometric convergence of backgrounds into the uniform comparison of functionals used above.

From this section on, each section specifies its own background and reference. Here we vary a Hermitian form $\beta$ on $X$ and use the zero reference. Write $B_\beta$, $H_\beta$ and $\mathcal A_\beta$ for the normalization scalar, its dual functional and the supporting sets, defined as above with $\theta=\beta$ and $\varphi=0$. For two such backgrounds, set
\[
 d(B_\beta,B_{\beta'})=\sup_{F\in L^\infty(\mathrm{d}V)}|B_\beta(F)-B_{\beta'}(F)|.
\]
The supremum ranges over all bounded measurable tests, with no common bound on their norms.

\begin{theorem}\label{thm:geometric-normalization-comparison}\index{Stability!normalization functional}
Suppose $0<\varepsilon<1$ and $(1-\varepsilon)\beta\leq\beta'\leq(1+\varepsilon)\beta$. For every bounded measurable $F$,
\begin{equation}\label{eq:geometric-normalization-comparison}
 n\log(1-\varepsilon)\leq B_{\beta'}(F)-B_\beta(F)\leq n\log(1+\varepsilon).
\end{equation}
In particular, $d(B_\beta,B_{\beta'})\leq-n\log(1-\varepsilon)$. The same conclusion holds if the two inequalities between forms are satisfied by positive representatives $\beta+\ddc h$ and $\beta'+\ddc h'$, with $h,h'$ smooth real functions.
\end{theorem}
\begin{proof}
First let $F$ be smooth, and let $v,v'$ be smooth solutions in the two backgrounds. At a maximum point of $v'-(1+\varepsilon)v$,
\[
 0<\beta'+\ddc v'\leq(1+\varepsilon)(\beta+\ddc v).
\]
Taking determinants and canceling $\mathrm{e}^F\mathrm{d}V$ gives the upper bound. At a minimum point of $v'-(1-\varepsilon)v$, the reverse inequality between forms gives the lower bound. For bounded measurable $F$, choose smooth uniformly bounded approximations converging in $L^1$ and apply \cref{thm:normalization-bounded-lone} in each of the two backgrounds. Finally, the correspondence $v\mapsto v-h$ preserves both boundedness and the current, so $B_{\beta+\ddc h}=B_\beta$ whenever both backgrounds are positive. This proves the assertion about representatives.
\end{proof}

\begin{corollary}\label{cor:geometric-adjoint-stability}
Let $\beta_j$ be Hermitian forms with $(1-\varepsilon_j)\beta\leq\beta_j\leq(1+\varepsilon_j)\beta$, where $0<\varepsilon_j<1$ and $\varepsilon_j\to0$. Suppose that $F_j\to F$ in $L^1(\mathrm{d}V)$ with a common $L^\infty$ bound. Then $B_{\beta_j}(F_j)\to B_\beta(F)$. For any $\nu_j\in\mathcal A_{\beta_j}(F_j)$, the densities are uniformly integrable, and every subsequence has a further subsequence converging weakly in $L^1$ to an element of $\mathcal A_\beta(F)$. If this set is a singleton, which is the case in particular when $F$ is smooth, the whole sequence converges to it weakly in $L^1$ and as probability measures.
\end{corollary}
\begin{proof}
Put $\delta_j=-n\log(1-\varepsilon_j)$. The scalar assertion follows from \eqref{eq:geometric-normalization-comparison} and \cref{thm:normalization-bounded-lone} in the fixed background $\beta$. Taking suprema over tests gives the inequalities, valid also for infinite values,
\[
 H_{\beta_j}\leq H_\beta+\delta_j,\qquad H_\beta\leq H_{\beta_j}+\delta_j.
\]
In particular, the two duals are finite on the same measures, and the identity $H_{\beta_j}(\nu_j)=B_{\beta_j}(F_j)+\int_XF_j\,\mathrm{d}\nu_j$ gives
\[
 0\leq H_\beta(\nu_j)-B_\beta(F_j)-\int_XF_j\,\mathrm{d}\nu_j\leq2\delta_j.
\]
If $\|F_j\|_\infty\leq M$, this puts the measures in the fixed sublevel set $\{H_\beta\leq B_\beta(0)+2M+2\sup_j\delta_j\}$, to which \cref{thm:normalization-dual-compactness} applies. For a weak $L^1$ cluster point, uniform integrability gives $\int_X|F_j-F|\,\mathrm{d}\nu_j\to0$, and lower semicontinuity in the last inequality, together with the minimum formula of \cref{thm:normalization-measurable-adjoint}, identifies the limit as a supporting probability at $F$. Compactness gives convergence of the whole sequence when that supporting probability is unique, and \cref{thm:normalization-gauduchon-recovery} gives uniqueness for smooth $F$.
\end{proof}

Smooth changes of the background also control the potentials, because the inverse of the linearized equation used in \cref{prop:smooth-dependence} remains available when the form itself is included among the parameters.

\begin{proposition}\label{prop:smooth-background-potential-stability}
Let $\beta_j,\beta$ be Hermitian forms and $F_j,F$ smooth real functions with $\beta_j\to\beta$ and $F_j\to F$ in $C^\infty$. Then the solutions $v_j$ for $(\beta_j,F_j)$ normalized by $\int_Xv_j\,\mathrm{d}V=0$ converge in $C^\infty$ to the similarly normalized solution $v$ for $(\beta,F)$, and their scalars converge as well.
\end{proposition}
\begin{proof}
Fix $k\geq0$ and $0<\alpha<1$. On the open set where $\eta+\ddc w>0$, consider the smooth map
\[
 (\eta,G,w,a)\longmapsto\log\frac{(\eta+\ddc w)^n}{\mathrm{d}V}-a-G
\]
defined for $C^{k,\alpha}$ forms $\eta$ and functions $G$, mean-zero $C^{k+2,\alpha}$ potentials $w$ and $a\in\mathbb R$, with values in $C^{k,\alpha}$. At the limiting solution its derivative in $(w,a)$ is $(h,q)\mapsto\Delta_\Omega h-q$, where $\Omega=\beta+\ddc v$ and $\Delta_\Omega h=n\ddc h\wedge\Omega^{n-1}/\Omega^n$. Given a target $g$, choose $q=-\int_Xg\,\mathrm{d}\mu_v$, where $\mu_v$ is the adjoint probability of \cref{prop:adjoint-gauduchon}. The Fredholm alternative then gives a unique mean-zero $h$ with $\Delta_\Omega h=g+q$, and the inverse is bounded as in \cref{prop:smooth-dependence}. In dimension one this is the Fredholm alternative for the Poisson equation.

The implicit function theorem gives a local branch of solutions depending continuously on the parameters in these spaces. For smooth parameters, uniqueness of bounded solutions modulo constants, \cref{cor:bounded-positive-density-uniqueness}, identifies this branch with the global smooth solution. Thus $v_j\to v$ in $C^{k+2,\alpha}$ for every $k$, which proves the assertion.
\end{proof}

The smooth adjoint probabilities determine the functional up to its value at zero. If $\nu_{\beta,F}$ denotes the unique supporting probability at a smooth $F$, integrating the derivative of \cref{thm:normalization-gauduchon-recovery} along the path $s\mapsto sF$ gives
\begin{equation}\label{eq:adjoint-functional-reconstruction}
 B_\beta(F)=B_\beta(0)-\int_0^1\!\int_XF\,\mathrm{d}\nu_{\beta,sF}\,\mathrm{d}s.
\end{equation}
Here the derivative is continuous along the path by the implicit function theorem. Consequently, two Hermitian backgrounds on the same $X$, with the same $\mathrm{d}V$, have the same functional on bounded measurable tests if and only if their values at zero and all their adjoint probabilities at smooth tests agree. One direction follows by differentiation; the other follows from \eqref{eq:adjoint-functional-reconstruction} for smooth tests and from bounded $L^1$ approximation for the remaining tests.
